% ============================================================
%  CONCLUSION GÉNÉRALE
% ============================================================
\chapter*{Conclusion générale et perspectives}
\addcontentsline{toc}{chapter}{Conclusion générale et perspectives}
\markboth{Conclusion générale}{Conclusion générale}

\hspace{1cm}Ce stage de Projet de Fin d'Année, réalisé sur sept semaines (15~juillet -- 31~août~2026) au sein du \gls{cmrpi} -- Espace Maroc Cyberconfiance, a porté sur la conception et le développement d'une plateforme d'auto-évaluation de la maturité cybersécurité destinée aux \gls{pme} marocaines. Au terme de ce travail, je peux dresser le bilan suivant.

\vspace{0.5cm}
\textbf{Bilan des réalisations :}

\begin{itemize}
  \item Une analyse des référentiels normatifs (guide \gls{cmrpi}/\gls{ausim}, \gls{iso27001}, \gls{nistcsf}, loi 09-08) a permis de justifier académiquement une structure de questionnaire en 5 domaines et 25 questions.
  \item Une plateforme fonctionnelle a été conçue et développée, intégrant un questionnaire d'auto-évaluation, un moteur de scoring automatisé et reproductible, ainsi qu'un export de rapports personnalisés au format \gls{pdf}.
  \item La qualité logicielle de la plateforme a été validée par 52 tests automatisés, 90\,\% de couverture de code et un pipeline d'intégration continue, avec une validation de la cohérence méthodologique sur trois profils fictifs de \gls{pme}.
  \item La démarche a été menée en trois jalons itératifs, avec un suivi hebdomadaire et une validation systématique par l'encadrant professionnel, le Dr.~Rachid Abouettahir.
\end{itemize}

\vspace{0.5cm}
\textbf{Compétences mobilisées :} analyse de référentiels normatifs en cybersécurité, conception d'une méthodologie de scoring justifiée académiquement, développement full-stack en Python, mise en place de tests automatisés et d'un pipeline d'intégration continue, gestion de projet et communication avec un encadrant professionnel.

\vspace{0.5cm}
\textbf{Perspectives :}

\begin{enumerate}
  \item Migrer l'architecture vers Django ou FastAPI côté serveur avec un frontend React, pour une version production multi-utilisateurs.
  \item Migrer la base de données de SQLite vers PostgreSQL afin de supporter une montée en charge.
  \item Déployer la plateforme sur Streamlit Community Cloud et finaliser la documentation technique.
  \item Mettre en place un benchmarking sectoriel anonymisé entre \gls{pme} et étudier la faisabilité d'une certification \enquote{PME Cyber-Sûre}.
\end{enumerate}

\vspace{0.5cm}
Au-delà du livrable technique, ce stage illustre l'impact potentiel d'une telle plateforme pour sensibiliser les \gls{pme} marocaines à leur maturité cybersécurité, en cohérence directe avec la mission du \gls{cmrpi}.
